%% =====================================================================
%% Appendices
%% =====================================================================

\section{Reproduction Guide}
\label{app:repro}

\paragraph{Prerequisites.} Docker Desktop (or Docker Engine with the Compose plugin),
about 2~GB of free RAM, and a free TCP port 8090 (configurable as \texttt{PUBLIC\_PORT}
in \texttt{.env}). Python~3 is needed only to run the test suite.

\paragraph{Start the stack.}
\begin{lstlisting}[language=bash, numbers=none]
docker compose up --build        # first run downloads images and builds (minutes)
docker compose ps                # all 9 services Up; db (healthy)
\end{lstlisting}

\paragraph{Endpoints.}
\begin{center}
\small
\begin{tabular}{@{}ll@{}}
\toprule
Dashboard (live operations console) & \url{http://localhost:8090} \\
CacheShop storefront & \url{http://localhost:8090/shop.html} \\
Product API & \url{http://localhost:8090/api/products/<id>} \\
Grafana & \url{http://localhost:3000} \\
Prometheus & \url{http://localhost:9090} \\
\bottomrule
\end{tabular}
\end{center}

\paragraph{Run the three experiments.} Use the dashboard presets \emph{Exp~1}
(baseline), \emph{Exp~2} (attack, no defense), and \emph{Exp~3} (attack with defenses),
or call the control API directly:
\begin{lstlisting}[language=bash, numbers=none]
curl -X POST localhost:8090/control/preset -H 'Content-Type: application/json' -d '{"exp":1}'
curl -X POST localhost:8090/control/preset -H 'Content-Type: application/json' -d '{"exp":2}'
curl -X POST localhost:8090/control/reset
curl -X POST localhost:8090/control/preset -H 'Content-Type: application/json' -d '{"exp":3}'
\end{lstlisting}
Allow at least 5--10~s of baseline traffic before an attack so that the ML detector can
train. Each monitor start creates a new \texttt{results/metrics\_<timestamp>.csv}.

\paragraph{Validate.}
\begin{lstlisting}[language=bash, numbers=none]
python tests/test_cachecollapse.py   # expected: SUMMARY: 16 passed, 0 failed
\end{lstlisting}

\paragraph{Stop.} \texttt{docker compose down} (add \texttt{-v} to remove volumes).

\section{Key Configuration Listings}
\label{app:listings}

\begin{lstlisting}[caption={Nginx: edge rate-limit zone, client-address forwarding, and the live-reloadable include (excerpt of \texttt{nginx/nginx.conf}).}, label={lst:nginx}]
limit_req_zone $binary_remote_addr zone=api_rl:10m rate=50r/s;
upstream backend { server backend:5000; keepalive 64; }

location /api/ {
    proxy_set_header X-Forwarded-For $remote_addr;   # overwrite, not append
    proxy_set_header X-Real-IP $remote_addr;
    include /etc/nginx/dynamic/ratelimit.conf;      # rewritten by the monitor
    proxy_http_version 1.1;
    proxy_set_header Connection "";
    proxy_pass http://backend/;
}
# dynamic/ratelimit.conf when ON:
#   limit_req zone=api_rl burst=20 nodelay;
\end{lstlisting}

\begin{lstlisting}[caption={Nginx container command: start Nginx and reload it whenever the edge rate-limit include changes (excerpt of \texttt{docker-compose.yml}).}, label={lst:watcher}]
nginx -g 'daemon off;' &
prev=""
while :; do
  sleep 1
  cur=$(cat /etc/nginx/dynamic/ratelimit.conf 2>/dev/null)
  if [ "$cur" != "$prev" ]; then prev="$cur"; nginx -s reload 2>/dev/null; fi
done
\end{lstlisting}

\begin{lstlisting}[language=Python, caption={Rule-based detector (\texttt{monitor/detect\_rules.py}).}, label={lst:rule}]
DB_MULTIPLIER = 5       # "high DB load" = more than 5x the learned baseline
BASELINE_FLOOR = 5      # never treat less than this many q/s as the baseline
RESP_MS_LIMIT = 200     # "high response time"
MISS_RATE_LIMIT = 0.5   # "high miss rate"

def evaluate(sample, baseline_db):
    base = max(baseline_db, BASELINE_FLOOR)
    c1 = sample["db_qps"] > DB_MULTIPLIER * base
    c2 = sample["resp_avg_ms"] > RESP_MS_LIMIT
    c3 = sample["miss_rate"] > MISS_RATE_LIMIT
    return {"db_load": c1, "response": c2, "miss_rate": c3}, (c1 and c2 and c3)
\end{lstlisting}

\begin{lstlisting}[language=Python, caption={IsolationForest detector core (\texttt{monitor/detect\_ml.py}).}, label={lst:ml}]
FEATURES = ["db_qps", "resp_avg_ms", "miss_rate", "rps", "failed"]
MIN_TRAIN = 10

def observe_baseline(self, sample):
    self._baseline.append([sample[f] for f in FEATURES])
    if len(self._baseline) >= MIN_TRAIN:
        self._baseline = self._baseline[-300:]
        data = np.array(self._baseline, dtype=float)
        self._scaler = StandardScaler().fit(data)
        self._model = IsolationForest(contamination=0.01, random_state=0)
        self._model.fit(self._scaler.transform(data))

def score(self, sample):
    if self._model is None:
        return 0.08, False
    x = self._scaler.transform([[sample[f] for f in FEATURES]])
    raw = float(self._model.decision_function(x)[0])  # +normal / -anomaly
    return max(0.0, min(1.0, 0.5 - raw * 4.0)), self._model.predict(x)[0] == -1
\end{lstlisting}

\section{Mapping to the Attributes of a Complex Engineering Problem}
\label{app:cep}
Table~\ref{tab:cep} maps the project onto the attributes of a complex engineering
problem defined in the International Engineering Alliance's graduate attributes
\cite{ieagapc2021}, which underlie the accreditation criteria applied to this course.

\begin{table}[h]
\centering
\caption{Complex engineering problem attributes addressed by CacheCollapse.}
\label{tab:cep}
\footnotesize
\begin{tabularx}{\linewidth}{@{}L{3.6cm}Y@{}}
\toprule
Attribute & How it is addressed \\
\midrule
WP1 Depth of knowledge & Caching theory, queueing and overload (Sec.~\ref{sec:bg-queueing}), rate-limiting algorithms, anomaly detection (Sec.~\ref{sec:bg-iforest}), container orchestration \\
WP2 Range of conflicting requirements & Database protection vs.\ user latency vs.\ data freshness vs.\ false alarms (tensions T1--T3, Sec.~\ref{sec:intro}) \\
WP3 Depth of analysis & Analytical model (Sec.~\ref{sec:model}), statistical evaluation, threshold sensitivity, ROC, ablation, longitudinal analysis (Sec.~\ref{sec:results}) \\
WP4 Familiarity of issues & Reframes a known performance hazard as an adversarial, emergent availability attack that cannot be patched \\
WP5 Extent of applicable codes & HTTP semantics (RFC 6585, RFC 5861), OWASP API Top 10, MITRE ATT\&CK, NIST SP 800-61 (Sec.~\ref{sec:threat}) \\
WP6 Stakeholder involvement & End users, operators, security analysts, instructor/audience; separate harm and novelty alerting for operators vs.\ analysts \\
WP7 Interdependence & Nine interacting services; defenses at different layers that are only effective in combination (Sec.~\ref{sec:ablation}) \\
\bottomrule
\end{tabularx}
\end{table}

\section{Glossary}
\label{app:glossary}
\begin{description}[leftmargin=1.2em,style=nextline,itemsep=1pt]
\item[Cache-aside] Pattern in which the application reads the cache and, on a miss,
loads from the database and fills the cache.
\item[Cache stampede (dog-pile, thundering herd)] Many concurrent misses for the same
key regenerating it simultaneously.
\item[Coalescing ratio $\rho$] Cache misses per database attempt; 1 without coalescing.
\item[Edge rate limiting] Limiting applied at the outermost proxy, before requests
consume upstream resources.
\item[Fail-fast pool] A connection pool that rejects requests immediately when all
connections are in use.
\item[Goodput] Useful completed work per unit time, as opposed to attempted work.
\item[IDDoS] Indirect (distributed) denial of service: overloading one component by
sending legitimate requests to another.
\item[IsolationForest] Unsupervised anomaly detector that scores points by how quickly
random trees isolate them.
\item[Single-flight] Allowing only one in-progress regeneration per key.
\item[Stale-while-revalidate (SWR)] Serving a slightly old copy while one request
refreshes the entry.
\item[TTL] Time-to-live of a cache entry.
\end{description}
